% !TeX root = ../main.tex
%% =======================================================================
%  RESUMEN + ABSTRACT


\chapter*{Resumen}
\addcontentsline{toc}{chapter}{Resumen}
\markboth{Resumen}{Resumen}

Hay un amplio consenso en que las comunicaciones \'opticas intersatelitales
representan una tecnolog\'ia clave para las futuras redes espaciales de alta
capacidad. En particular, los enlaces entre sat\'elites de \'orbita baja (LEO)
y geoestacionaria (GEO) ofrecen una alternativa para transmitir datos desde
constelaciones LEO hacia estaciones terrestres, reduciendo la dependencia de la
visibilidad directa entre los sat\'elites y la Tierra. Sin embargo, las tasas
de transmisi\'on de los enlaces \'opticos LEO--GEO actualmente operativos
permanecen en torno a \SI{1.8}{\giga\bit\per\second}, mientras que las
comunicaciones entre sat\'elites LEO ya alcanzan aproximadamente
\SI{100}{\giga\bit\per\second}. Esta diferencia motiva el desarrollo de
sistemas LEO--GEO de mayor capacidad, basados en t\'ecnicas avanzadas de
modulaci\'on, multiplexaci\'on en polarizaci\'on y recepci\'on coherente con
procesamiento digital de se\~nales.

Se diseñ\'o y verific\'o mediante simulaciones numéricas e implementación experimental un receptor \'optico
coherente con modulaci\'on QPSK en doble polarizaci\'on (DP-QPSK) de \SI{128}{\giga\bit\per\second} (\SI{32}{\giga\baud}) para un enlace LEO--GEO a
\SI{1550}{\nano\metre}. Se model\'o la cadena completa, con el transmisor, el
canal satelital con Doppler y rotaci\'on de polarizaci\'on, y el receptor
coherente. Se implementaron los algoritmos del DSP para la
ortogonalizaci\'on, la recuperaci\'on de frecuencia y de reloj, la ecualizaci\'on
adaptativa MIMO $2\times2$ y la recuperaci\'on de fase.
Se mont\'o la cadena en el laboratorio y se verific\'o el correcto
funcionamiento del DSP para tres tasas de transmisi\'on:
\SI{12.5}{\giga\baud}, \SI{25}{\giga\baud} y \SI{32}{\giga\baud}.

El receptor cumple los requisitos del enlace. Para la OSNR objetivo y valores
realistas de los componentes, el balance de potencia cierra con \SI{3.4}{\dB}
de margen. La tolerancia al desplazamiento de frecuencia por efecto Doppler
sigue la relaci\'on $\abs{\Df} < \Rs/4$, verificada experimentalmente en las tres
tasas. A la tasa de dise\~no equivale a \SI{8.0}{\giga\hertz} de tolerancia, contra los 
\SI{6.0}{\giga\hertz} de requisito que impone el enlace. La BER medida en el laboratorio a
esta tasa es de \num{1.2e-3}, por debajo del umbral HD-FEC. Incluyendo la
penalidad de implementaci\'on del receptor y los deterioros esperados en vuelo,
el umbral se alcanza con una OSNR de \SI{13.6}{\dB}. Frente a los dem\'as
deterioros del canal, la rotaci\'on de polarizaci\'on, el chirp Doppler y el
ancho de l\'inea de los l\'aseres, las tolerancias medidas del receptor superan
entre uno y diez \'ordenes de magnitud los valores que el enlace presenta.

La arquitectura propuesta resulta entonces viable para enlaces \'opticos
LEO--GEO de alta capacidad. El margen del enlace no queda limitado por el
procesamiento digital sino por el balance de potencia y el apuntamiento.

\vspace{1.5em}
\noindent\textbf{Palabras clave:}
\textsc{comunicaciones \'opticas coherentes, dp-qpsk, enlaces intersatelitales,
procesamiento digital de se\~nales, efecto doppler}

\cleardoublepage

\chapter*{Abstract}
\addcontentsline{toc}{chapter}{Abstract}
\markboth{Abstract}{Abstract}

\selectlanguage{english}
\sisetup{output-decimal-marker = {.}}
There is broad agreement that optical inter-satellite links are a key
technology for future high-capacity space networks. In particular, links
between low Earth orbit (LEO) and geostationary orbit (GEO) satellites offer a
way to relay data from LEO constellations to ground stations, reducing the
dependence on direct visibility between the satellites and the Earth. However,
the transmission rates of the LEO--GEO optical links currently in operation
remain around \SI{1.8}{\giga\bit\per\second}, while communications between LEO
satellites already reach about \SI{100}{\giga\bit\per\second}. This gap
motivates the development of higher-capacity LEO--GEO systems based on advanced
modulation, polarization multiplexing and coherent reception with digital
signal processing.

A coherent optical receiver using dual-polarization quadrature phase-shift
keying (DP-QPSK) at \SI{128}{\giga\bit\per\second} (\SI{32}{\giga\baud}) for a
LEO--GEO link at \SI{1550}{\nano\metre} was designed and verified through. The
complete chain was modelled, including the transmitter, the satellite channel
with Doppler and polarization rotation, and the coherent receiver. The DSP
algorithms were implemented for orthogonalization, frequency and clock
recovery, adaptive $2\times2$ MIMO equalization and carrier phase recovery. The
chain was built in the laboratory and the DSP was verified at three symbol
rates, \SI{12.5}{\giga\baud}, \SI{25}{\giga\baud} and \SI{32}{\giga\baud}.

The receiver meets the link requirements. For the target OSNR and realistic
component parameters, the power budget closes with \SI{3.4}{\dB} of margin. The
tolerance to Doppler frequency offset follows $\abs{\Df} < \Rs/4$, verified
experimentally at the three symbol rates. At the design rate this limit amounts
to \SI{8.0}{\giga\hertz}, against the \SI{6.0}{\giga\hertz} required by the
link. The BER measured in the laboratory at that rate is \num{1.2e-3}, below
the HD-FEC threshold. Including the implementation penalty of the receiver and
the impairments expected in flight, that threshold is met at an OSNR of
\SI{13.6}{\dB}. For the remaining channel impairments, polarization rotation,
Doppler chirp and laser linewidth, the measured tolerances of the receiver
exceed the values imposed by the link by one to ten orders of magnitude.

The proposed architecture is therefore viable for high-capacity LEO--GEO
optical links. The link margin is set not by the digital signal processing but
by the power budget and the pointing accuracy.

\vspace{1.5em}
\noindent\textbf{Keywords:}
\textsc{coherent optical communications, dp-qpsk, inter-satellite links, digital
signal processing, doppler effect}

\sisetup{output-decimal-marker = {,}}
\selectlanguage{spanish}


\cleardoublepage
